\chapter{पिरामिड और शंकु}\label{ch:g8:solids}

\cref{ch:g7:areas} वाले \omterm{def:g7:areas:prism}{प्रिज़्म} और \omterm{def:g7:areas:prism}{बेलन} --- यानी हर जगह एक जैसी
काट वाले \omterm{def:g5:solids:def}{ठोस} --- के बाद अब आते हैं नोक वाले \omterm{def:g5:solids:def}{ठोस}: \omterm{def:g8:solids:def}{पिरामिड} और \omterm{def:g8:solids:def}{शंकु}।
उनके \omterm{def:g6:measure:volume}{आयतन} के सूत्र में एक मशहूर गुणक $\frac13$ बैठा है, जिसे यह
अध्याय मानने लायक़ बनाता है और इस्तेमाल करता है; \omterm{def:g5:solids:def}{ठोसों} की पढ़ाई
\cref{ch:g9:solids} में गोलों और काटों के साथ आगे बढ़ती है।

\section{पिरामिड और शंकु का ब्योरा}

\begin{definition}[पिरामिड, शंकु]\label{def:g8:solids:def}
\begin{itemize}
  \item \emph{पिरामिड}\index{पिरामिड} का \emph{आधार} एक \omterm{def:g4:geometry:polygon}{बहुभुज} होता
        है और एक बिंदु, \emph{शिखर}, जो आधार के हर \omterm{def:g5:solids:def}{शीर्ष} से
        त्रिभुजाकार फलकों से जुड़ा होता है। उसकी \emph{ऊँचाई} शिखर
        से आधार के तल तक की दूरी है।
  \item \emph{शंकु}\index{शंकु} वही रचना है, बस चक्र के ऊपर: एक
        आधार-चक्र, एक शिखर, और एक घुमावदार बग़ली सतह।
\end{itemize}
पिरामिड \emph{सम} तब कहलाता है जब उसका आधार एक सम \omterm{def:g4:geometry:polygon}{बहुभुज} हो (वर्ग,
\omterm{def:g6:shapes:triangles}{समबाहु त्रिभुज}, \dots) और उसका शिखर आधार के केंद्र के ठीक ऊपर बैठा
हो।
\end{definition}

\begin{omfigure}
\begin{tikzpicture}[scale=0.75]
  % square pyramid
  \coordinate (P1) at (-1.3,0);
  \coordinate (P2) at (0.9,0);
  \coordinate (P3) at (1.7,0.75);
  \coordinate (P4) at (-0.5,0.75);
  \coordinate (S) at (0.2,3.0);
  \draw[thick] (P1) -- (P2) -- (P3);
  \draw[dashed] (P3) -- (P4) -- (P1);
  \draw[thick] (P1) -- (S) -- (P2);
  \draw[thick] (P3) -- (S);
  \draw[dashed] (P4) -- (S);
  \draw[dashed] (S) -- (0.2,0.375);
  \node[right, font=\small] at (0.22,1.6) {$h$};
  \node[above, font=\small] at (S) {शिखर};
  \node[below, font=\small] at (-0.2,-0.1) {वर्गाकार आधार};
  % cone
  \begin{scope}[xshift=6.2cm]
  \draw[thick] (0,0) ellipse (1.5 and 0.5);
  \coordinate (T) at (0,3.0);
  \draw[thick] (-1.5,0) -- (T) -- (1.5,0);
  \draw[dashed] (T) -- (0,0);
  \node[right, font=\small] at (0.02,1.5) {$h$};
  \draw[thin] (0,0) -- (1.5,0) node[midway, below, font=\small] {$r$};
  \end{scope}
\end{tikzpicture}

{\small वर्गाकार आधार वाला एक \omterm{def:g8:solids:def}{पिरामिड} और एक \omterm{def:g8:solids:def}{शंकु}: एक बहुभुजाकार या
वृत्ताकार आधार, एक शिखर, और आधार पर \omterm{def:g6:lines:perp}{लंब} नापी गई ऊँचाई।}
\end{omfigure}

\begin{example}[फलक, किनारे और शीर्ष गिनना]\label{ex:g8:solids:count}
वर्गाकार आधार वाले \omterm{def:g8:solids:def}{पिरामिड} के $5$ \omterm{def:g5:solids:def}{फलक} ($1$ वर्ग $+ 4$ त्रिभुज),
$8$ \omterm{def:g5:solids:def}{किनारे} और $5$ \omterm{def:g5:solids:def}{शीर्ष} होते हैं। $n$ भुजाओं वाले आधार पर:
$n + 1$ \omterm{def:g5:solids:def}{फलक}, $2n$ \omterm{def:g5:solids:def}{किनारे}, $n + 1$ \omterm{def:g5:solids:def}{शीर्ष}। ऑयलर का नन्हा ढर्रा जाँचो:
\omterm{def:g5:solids:def}{फलक} $+$ \omterm{def:g5:solids:def}{शीर्ष} $=$ \omterm{def:g5:solids:def}{किनारे} $+ 2$।
\end{example}

\section{आयतन}

\begin{theorem}[पिरामिड या शंकु का आयतन]\label{thm:g8:solids:volume}
आधार का \omterm{def:g6:measure:area}{क्षेत्रफल} $B$ और ऊँचाई $h$ वाले \omterm{def:g8:solids:def}{पिरामिड} या \omterm{def:g8:solids:def}{शंकु} के लिए:
\[
V = \frac{1}{3}\, B \times h
\]
--- यानी उसी आधार और उसी ऊँचाई वाले \omterm{def:g7:areas:prism}{प्रिज़्म} या \omterm{def:g7:areas:prism}{बेलन} का एक तिहाई।
आधार-त्रिज्या $r$ वाले \omterm{def:g8:solids:def}{शंकु} के लिए: $V = \frac13 \pi r^2 h$।
\end{theorem}
\admitted

\begin{remark}[एक तिहाई क्यों?]
किसी खोखले \omterm{def:g8:solids:def}{पिरामिड} को पानी से भरो और उसी आधार तथा ऊँचाई वाले
\omterm{def:g7:areas:prism}{प्रिज़्म} में उँडेलो: ठीक तीन बार भरना पड़ता है। और भी अच्छी बात: किसी
\omterm{def:g5:solids:def}{घन} को तीन एक जैसे पिरामिडों में काटा जा सकता है, और हर एक का आधार \omterm{def:g5:solids:def}{घन}
का कोई \omterm{def:g5:solids:def}{फलक} होता है तथा शिखर \omterm{def:g5:solids:def}{घन} का कोई \omterm{def:g5:solids:def}{शीर्ष} (मन में चित्र बनाने की
कोशिश करो!)। आम प्रमाण समाकलन माँगता है, जो हाई-स्कूल खंड का औज़ार है।
\end{remark}

\begin{example}\label{ex:g8:solids:volume}
किसी \omterm{def:g8:solids:def}{पिरामिड} का वर्गाकार आधार $5$ cm भुजा का है और ऊँचाई $9$ cm:
\begin{enumerate}
  \item आधार का \omterm{def:g6:measure:area}{क्षेत्रफल}: $B = 5^2 = 25$ cm$^2$;
  \item \omterm{def:g6:measure:volume}{आयतन}: $V = \frac13 \times 25 \times 9 = 25 \times 3 = 75$
        cm$^3$।
\end{enumerate}
$3$ cm \omterm{def:g3:shapes:shapes}{त्रिज्या} और $10$ cm ऊँचाई वाला एक कुल्फ़ी का \omterm{def:g8:solids:def}{शंकु}:
\[
V = \frac13 \times \pi \times 3^2 \times 10 = 30\pi \approx 94
\text{ cm}^3 .
\]
\end{example}

\begin{example}[ऊँचाई से सावधान]\label{ex:g8:solids:slant}
\omterm{def:g8:solids:def}{शंकु} में ऊँचाई $h$ (शिखर से आधार के केंद्र तक, \omterm{def:g6:lines:perp}{लंब}) को \emph{तिरछी
ऊँचाई} $s$ (शिखर से किनार तक) से मत उलझाओ। दोनों शिखर--केंद्र--किनार
वाले \omterm{def:g6:shapes:triangles}{समकोण त्रिभुज} में पाइथागोरस
(\cref{thm:g8:pythagoras:direct}) से जुड़े हैं:
\[
s^2 = h^2 + r^2 .
\]
इसलिए जिस \omterm{def:g8:solids:def}{शंकु} में $r = 3$ और $s = 5$ हो उसकी ऊँचाई
$h = \sqrt{25 - 9} = 4$ है, और \omterm{def:g6:measure:volume}{आयतन}
$\frac13 \pi \times 9 \times 4 = 12\pi$।
\end{example}

\begin{method}[आयतन वाली समस्याएँ]\label{met:g8:solids:method}
\begin{enumerate}
  \item \omterm{def:g5:solids:def}{ठोस} पहचानो (\omterm{def:g7:areas:prism}{प्रिज़्म}/\omterm{def:g7:areas:prism}{बेलन}: $V = Bh$; \omterm{def:g8:solids:def}{पिरामिड}/\omterm{def:g8:solids:def}{शंकु}:
        $V = \frac13 Bh$);
  \item पहले, अलग क़दम के तौर पर, आधार का \omterm{def:g6:measure:area}{क्षेत्रफल} $B$ निकालो;
  \item जाँचो कि ऊँचाई आधार पर \omterm{def:g6:lines:perp}{लंब} है (तिरछी ऊँचाई दी हो तो
        पाइथागोरस लगाओ);
  \item गुणा करो, मात्रक एक जैसे रखो, और ज़रूरत हो तो आख़िर में बदल
        लो ($1$ L $= 1000$ cm$^3$)।
\end{enumerate}
\end{method}

\section{जाल}

\begin{example}[पिरामिड का जाल]\label{ex:g8:solids:net}
वर्गाकार आधार वाले सम \omterm{def:g8:solids:def}{पिरामिड} को खोलने पर वह चपटा होकर अपना
\emph{जाल} बन जाता है: बीच में वर्गाकार आधार और उसकी भुजाओं से जुड़े
चार एक जैसे \omterm{def:g6:shapes:triangles}{समद्विबाहु त्रिभुज}। इन त्रिभुजों की बराबर भुजाएँ \omterm{def:g8:solids:def}{पिरामिड}
के \emph{बग़ली \omterm{def:g5:solids:def}{किनारे}} हैं --- और उनकी लंबाई न $h$ है और न किसी \omterm{def:g5:solids:def}{फलक}
की तिरछी ऊँचाई, इसलिए गत्ता काटने से पहले नाप ध्यान से लिख लो।
\end{example}

\begin{omfigure}
\begin{tikzpicture}[scale=0.62]
  \draw[thick] (0,0) rectangle (3,3);
  \draw[thick] (0,0) -- (1.5,-2.2) -- (3,0);
  \draw[thick] (3,0) -- (5.2,1.5) -- (3,3);
  \draw[thick] (0,3) -- (1.5,5.2) -- (3,3);
  \draw[thick] (0,0) -- (-2.2,1.5) -- (0,3);
  \node[font=\small] at (1.5,1.5) {आधार};
\end{tikzpicture}

{\small वर्गाकार आधार वाले \omterm{def:g8:solids:def}{पिरामिड} का जाल: चारों त्रिभुजों को तब तक
ऊपर मोड़ो जब तक उनके शिखर आपस में न मिल जाएँ।}
\end{omfigure}

\section{अभ्यास}

\begin{exercise}[$\star$]\label{exo:g8:solids:1}
त्रिभुजाकार आधार वाले \omterm{def:g8:solids:def}{पिरामिड} (\emph{चतुष्फलक}) के कितने \omterm{def:g5:solids:def}{फलक},
\omterm{def:g5:solids:def}{किनारे} और \omterm{def:g5:solids:def}{शीर्ष} होते हैं? और षट्भुजाकार आधार वाले के?
\end{exercise}

\begin{exercise}[$\star$]\label{exo:g8:solids:2}
$36$ cm$^2$ आधार-क्षेत्रफल और $10$ cm ऊँचाई वाले \omterm{def:g8:solids:def}{पिरामिड} का \omterm{def:g6:measure:volume}{आयतन}
निकालो; और $6 \times 4$ cm के आयताकार आधार तथा $7.5$ cm ऊँचाई वाले
\omterm{def:g8:solids:def}{पिरामिड} का।
\end{exercise}

\begin{exercise}[$\star$]\label{exo:g8:solids:3}
$5$ cm \omterm{def:g3:shapes:shapes}{त्रिज्या} और $9$ cm ऊँचाई वाले \omterm{def:g8:solids:def}{शंकु} का \omterm{def:g6:measure:volume}{आयतन} निकालो ($\pi$ के
साथ ठीक मान, फिर cm$^3$ तक गोल)।
\end{exercise}

\begin{exercise}[$\star$]\label{exo:g8:solids:4}
एक \omterm{def:g7:areas:prism}{बेलन} और एक \omterm{def:g8:solids:def}{शंकु} की \omterm{def:g3:shapes:shapes}{त्रिज्या} एक ही, $4$ cm, है और ऊँचाई भी एक ही,
$12$ cm। दोनों के \omterm{def:g6:measure:volume}{आयतन} निकालो। उनका अनुपात क्या है?
\end{exercise}

\begin{exercise}[$\star$]\label{exo:g8:solids:5}
किसी \omterm{def:g8:solids:def}{पिरामिड} का \omterm{def:g6:measure:volume}{आयतन} $56$ cm$^3$ और ऊँचाई $8$ cm है। उसके आधार का
\omterm{def:g6:measure:area}{क्षेत्रफल} क्या है? (\omterm{def:g6:measure:volume}{आयतन} का सूत्र लिखकर हल करो।)
\end{exercise}

\begin{exercise}[$\star$]\label{exo:g8:solids:6}
किसी \omterm{def:g8:solids:def}{शंकु} का जाल बनाओ (आधार के लिए एक चक्र, और बग़ली सतह के लिए किसी
बड़े चक्र की एक फाँक)। उस फाँक की \omterm{def:g3:shapes:shapes}{त्रिज्या} \omterm{def:g8:solids:def}{शंकु} की कौन-सी नाप है:
ऊँचाई $h$ या तिरछी ऊँचाई $s$?
\end{exercise}

\begin{exercise}[$\star\star$]\label{exo:g8:solids:7}
लूव्र के \omterm{def:g8:solids:def}{पिरामिड} का वर्गाकार आधार क़रीब $35$ m भुजा का है और ऊँचाई
क़रीब $22$ m। उसका \omterm{def:g6:measure:volume}{आयतन} नज़दीकी सौ \omterm{def:g5:solids:def}{घन} मीटर तक अंदाज़ो।
\end{exercise}

\begin{exercise}[$\star\star$]\label{exo:g8:solids:8}
किसी \omterm{def:g8:solids:def}{शंकु} की तिरछी ऊँचाई $13$ cm और आधार-त्रिज्या $5$ cm है। उसकी
ऊँचाई निकालो, फिर \omterm{def:g6:measure:volume}{आयतन} ($\pi$ के साथ ठीक मान)।
\end{exercise}

\begin{exercise}[$\star\star$]\label{exo:g8:solids:9}
$4$ cm \omterm{def:g3:shapes:shapes}{त्रिज्या} और $9$ cm ऊँचाई वाला एक शंक्वाकार गिलास \omterm{def:g5:solids:def}{किनारे} तक
भरा जाता है, फिर $4$ cm \omterm{def:g3:shapes:shapes}{त्रिज्या} वाले एक बेलनाकार गिलास में उँडेल
दिया जाता है। \omterm{def:g7:areas:prism}{बेलन} में तरल कितनी ऊँचाई तक पहुँचता है?
\end{exercise}

\begin{exercise}[$\star\star$]\label{exo:g8:solids:10}
एक बालू-घड़ी दो एक जैसे शंकुओं (\omterm{def:g3:shapes:shapes}{त्रिज्या} $2$ cm, ऊँचाई $4.5$ cm) से
बनी है, जो अपने शिखरों पर जुड़े हैं। सारी बालू ठीक एक \omterm{def:g8:solids:def}{शंकु} भर देती
है। बालू का \omterm{def:g6:measure:volume}{आयतन} निकालो (ठीक मान, फिर cm$^3$ में दशमलव के एक अंक तक
गोल), और यह भी कि बालू घड़ी के कुल भीतरी \omterm{def:g6:measure:volume}{आयतन} का कितना भाग घेरती है।
\end{exercise}

\begin{exercise}[$\star\star\star$]\label{exo:g8:solids:11}
$6$ cm आधार-भुजा वाले एक वर्गाकार-आधार \omterm{def:g8:solids:def}{पिरामिड} का शिखर आधार के एक
कोने के ठीक ऊपर, $8$ cm ऊँचाई पर है।
\begin{enumerate}
  \item क्या सूत्र $V = \frac13 Bh$ अब भी लागू होता है? (हाँ, होता
        है --- शिखर का बीचोंबीच होना ज़रूरी नहीं। $V$ निकालो।)
  \item शिखर से आधार के सामने वाले कोने तक जाते सबसे लंबे बग़ली
        \omterm{def:g5:solids:def}{किनारे} की लंबाई निकालो (पाइथागोरस दो बार: पहले आधार का
        विकर्ण)।
\end{enumerate}
\end{exercise}

\section{समस्या: तीन पिरामिडों में कटा हुआ घन}

\begin{problem}[{सप्ताहांत समस्या --- वह गुणक $\frac13$ आता कहाँ से
है, और आधी ऊँचाई पर कटे पिरामिड का क्या होता है}]
\label{pb:g8:solids:1}
\cref{thm:g8:solids:volume} के बाद वाली टिप्पणी दावा करती है कि किसी
\omterm{def:g5:solids:def}{घन} को तीन एक जैसे पिरामिडों में काटा जा सकता है --- ``मन में चित्र
बनाने की कोशिश करो!''। यह समस्या वह चित्र ठीक-ठीक बनाती है, उसी से
मशहूर गुणक $\frac13$ ईमानदारी से निकालती है, फिर उसी \omterm{def:g5:solids:def}{घन} को दूसरी तरह
\emph{छह} पिरामिडों में काटती है, और आख़िर में किसी \omterm{def:g8:solids:def}{पिरामिड} को उसकी
आधी ऊँचाई पर काटती है --- ऐसे नतीजे के साथ जो पहली बार में कम ही लोग
ठीक बता पाते हैं।

पूरे समय $ABCDEFGH$ भुजा $a$ का एक \omterm{def:g5:solids:def}{घन} है: आधार $ABCD$, ऊपरी \omterm{def:g5:solids:def}{फलक}
$EFGH$, जिसमें $E$, $A$ के ऊपर है, $F$, $B$ के ऊपर, $G$, $C$ के ऊपर
और $H$, $D$ के ऊपर।

\medskip
\noindent\textbf{भाग I --- एक ही शिखर वाले तीन \omterm{def:g8:solids:def}{पिरामिड}।}
\begin{enumerate}
  \item \omterm{def:g5:solids:def}{शीर्ष} $G$ \omterm{def:g5:solids:def}{घन} के तीन फलकों में पड़ता है। उन तीन फलकों की सूची
        बनाओ जिनमें $G$ \emph{नहीं} है, और उन तीन पिरामिडों का
        ब्योरा दो जो हर बार शिखर $G$ रखकर इनमें से एक-एक को आधार
        मानने पर बनते हैं। हर \omterm{def:g8:solids:def}{पिरामिड} के कितने \omterm{def:g5:solids:def}{फलक}, \omterm{def:g5:solids:def}{किनारे} और \omterm{def:g5:solids:def}{शीर्ष}
        हैं (\cref{ex:g8:solids:count})?
  \item समझाओ कि \omterm{def:g5:solids:def}{घन} को उसके लंबे विकर्ण $(AG)$ के चारों ओर एक तिहाई
        चक्कर घुमाने से \omterm{def:g5:solids:def}{घन} ज्यों का त्यों क्यों रहता है, पर सवाल 1
        के तीनों आधार आपस में चक्र में बदल जाते हैं --- और इसलिए
        तीनों \omterm{def:g8:solids:def}{पिरामिड} एक-दूसरे की एक जैसी नक़लें क्यों हैं।
  \item यह मानते हुए कि तीनों \omterm{def:g8:solids:def}{पिरामिड} \omterm{def:g5:solids:def}{घन} को बिना एक-दूसरे पर चढ़े भर
        देते हैं (गत्ते का नमूना काफ़ी क़ायल कर देता है), हर एक का
        \omterm{def:g6:measure:volume}{आयतन} निकालो।
  \item इसे \cref{thm:g8:solids:volume} के सूत्र से मिलाकर जाँचो:
        आधार $ABCD$ और शिखर $G$ वाले \omterm{def:g8:solids:def}{पिरामिड} का शिखर उसके आधार के
        एक \emph{कोने} के ठीक ऊपर है, ठीक वैसे ही जैसे
        \cref{exo:g8:solids:11} में। उसकी ऊँचाई पहचानो, सूत्र
        लगाओ, और तुलना करो।
  \item $6$ cm भुजा वाले \omterm{def:g5:solids:def}{घन} के लिए तीनों पिरामिडों में से हर एक का
        \omterm{def:g6:measure:volume}{आयतन} बताओ --- और समझाओ कि इस टुकड़ों में बाँटने का अध्याय
        वाले पानी उँडेलने के प्रयोग (एक \omterm{def:g7:areas:prism}{प्रिज़्म} के लिए \omterm{def:g8:solids:def}{पिरामिड} को
        तीन बार भरना) से क्या नाता है।
\end{enumerate}

\medskip
\noindent\textbf{भाग II --- केंद्र वाले छह \omterm{def:g8:solids:def}{पिरामिड}।}
अब मान लो $O$ \omterm{def:g5:solids:def}{घन} का केंद्र है, और $O$ को छहों फलकों में से हर एक के
चारों \omterm{def:g5:solids:def}{शीर्षों} से जोड़ो।
\begin{enumerate}[resume]
  \item इससे बनने वाले छह पिरामिडों का ब्योरा दो, और समझाओ कि वे
        एक-दूसरे की एक जैसी नक़लें क्यों हैं। हर एक की ऊँचाई क्या
        है?
  \item \omterm{def:g5:solids:def}{घन} का \omterm{def:g6:measure:volume}{आयतन} बाँटकर छहों पिरामिडों में से हर एक का \omterm{def:g6:measure:volume}{आयतन}
        निकालो।
  \item वही \omterm{def:g6:measure:volume}{आयतन} सूत्र $V = \frac13 B h$ से दोबारा निकालो, और
        जाँचो कि दोनों उत्तर मिलते हैं।
  \item ये दोनों बँटवारे कुछ बहुत ख़ास पिरामिडों के लिए गुणक
        $\frac13$ की पुष्टि करते हैं। समझाओ कि वे हर \omterm{def:g8:solids:def}{पिरामिड} और
        \omterm{def:g8:solids:def}{शंकु} के लिए \cref{thm:g8:solids:volume} अभी तक \emph{सिद्ध}
        क्यों नहीं करते --- और बताओ कि अध्याय इस कमी से कैसे
        निपटता है: कौन-सा कथन मान लिया गया है, कौन-सा प्रयोग उसका
        साथ देता है, और श्रृंखला के किस खंड में उसे ईमानदारी से
        सिद्ध किया जाता है।
  \item $6$ cm भुजा वाला \omterm{def:g5:solids:def}{घन} अपने छह केंद्र-पिरामिडों में काटा जाता
        है। हर एक का \omterm{def:g6:measure:volume}{आयतन} दोनों तरीक़ों से निकालो।
\end{enumerate}

\medskip
\noindent\textbf{भाग III --- \omterm{def:g8:solids:def}{पिरामिड} को आधी ऊँचाई पर काटना।}
$SABCD$ वर्गाकार आधार वाला एक सम \omterm{def:g8:solids:def}{पिरामिड} है: आधार-भुजा $c$, ऊँचाई
$h$, और शिखर $S$ आधार के केंद्र के ऊपर। उसे उस तल से काटो जो चारों
बग़ली \omterm{def:g5:solids:def}{किनारों} $[SA]$, $[SB]$, $[SC]$, $[SD]$ के \emph{मध्यबिंदुओं}
से गुज़रता है।
\begin{enumerate}[resume]
  \item त्रिभुज $SAB$ में मध्यबिंदु प्रमेय
        (\cref{thm:g8:midpoints:direct}) लगाओ: $[SA]$ और $[SB]$ के
        मध्यबिंदुओं को जोड़ते \omterm{def:g6:lines:objects}{रेखाखंड} की दिशा और लंबाई क्या है?
        इससे निकालो कि यह काट भुजा $\frac{c}{2}$ का एक वर्ग है, जो
        आधार के \omterm{def:g6:lines:perp}{समांतर} है और ऊँचाई $\frac{h}{2}$ पर बैठा है।
  \item काट के ऊपर वाला टुकड़ा ख़ुद एक वर्गाकार-आधार \omterm{def:g8:solids:def}{पिरामिड} है।
        उसकी आधार-भुजा और ऊँचाई बताओ, और दिखाओ कि उसका \omterm{def:g6:measure:volume}{आयतन} असली
        \omterm{def:g6:measure:volume}{आयतन} $V$ का ठीक \emph{एक बटा आठ} है।
  \item इससे नीचे वाले टुकड़े (\emph{छिन्नक}, यानी अधूरे \omterm{def:g8:solids:def}{पिरामिड} की
        शक्ल) का \omterm{def:g6:measure:volume}{आयतन} निकालो। आधी ऊँचाई वाली काट \omterm{def:g6:measure:volume}{आयतन} को किस
        अनुपात में बाँटती है?
  \item लूव्र का \omterm{def:g8:solids:def}{पिरामिड} (\cref{exo:g8:solids:7}: आधार-भुजा
        $35$ m, ऊँचाई $22$ m) दो चरणों में साफ़ किया जाता है: पहले
        आधी ऊँचाई से ऊपर का काँच, फिर नीचे का। आधी ऊँचाई से नीचे
        \emph{\omterm{def:g6:measure:volume}{आयतन}} का कितना भाग बैठता है? उसे \omterm{def:g5:solids:def}{घन} मीटर में, नज़दीकी
        सौ तक, अंदाज़ो।
  \item आम सीख: अगर किसी \omterm{def:g8:solids:def}{पिरामिड} की \emph{सारी} नापें $2$ से गुणा
        कर दी जाएँ तो उसके \omterm{def:g6:measure:volume}{आयतन} का क्या होता है? और हर नाप $k$ से
        गुणा करने पर \omterm{def:g6:measure:volume}{आयतन} कितने गुना बढ़ता है --- तथा इससे सवाल 12
        वाला ``एक बटा आठ'' एक पंक्ति में कैसे समझ आ जाता है?
\end{enumerate}
\end{problem}
